% TOP_PROBLEM: {"id": 3100083, "title": "What is the threshold function n = f(k) for the event that a random permutation on n symbols contains all patterns on k", "category_id": 2, "category": {"id": 2, "name": "combinatorics", "display_name": "Combinatorics", "description": "Counting problems, graph theory, discrete structures.", "slug": "combinatorics", "order_index": 2, "created_at": "2026-07-31T15:26:25.671Z"}, "status": "partially_solved", "problem_number": "AMR-030-0083", "set_id": 13, "statusQualification": "Corrects Cooper’s k²/2 constant to the standard k²/4 formulation in He–Kwan; the lower-threshold obstruction is already known."}
% FIRST_POSTED: 2026-10-02
% ENTRY_KIND: statement_only
% UnsolvedMath ID 3100083; source catalogue code AMR-030-0083
% !TeX program = lualatex
% Definition and sources only; this file is not an LLM solution attempt.
\documentclass[11pt,a4paper]{article}
\usepackage[margin=25mm]{geometry}
\usepackage{fontspec}
\setmainfont{lmroman10-regular.otf}[BoldFont=lmroman10-bold.otf,
 ItalicFont=lmroman10-italic.otf,BoldItalicFont=lmroman10-bolditalic.otf]
\usepackage{amsmath,amssymb,mathtools}
\usepackage{unicode-math}
\setmathfont{latinmodern-math.otf}
\AtBeginDocument{\let\setminus\smallsetminus}
\usepackage{xurl,hyperref}
\hypersetup{colorlinks=true,urlcolor=blue}
\setcounter{secnumdepth}{0}
\setlength{\parindent}{0pt}
\setlength{\parskip}{5pt}
\setlength{\emergencystretch}{3em}
\begin{document}
\section*{What is the threshold function n = f(k) for the event that a random permutation on n symbols contains all patterns on k}\label{3100083}
{\small UnsolvedMath ID 3100083 · AMR-030-0083 · Combinatorics · AMR Open Problem Lists}
\subsection{Definitions and mathematical statement}
For $\pi\in S_n$ and $\sigma\in S_k$, say that $\pi$ contains
$\sigma$ if some indices $i_1<\cdots<i_k$ have values in the
same relative order as $\sigma$. A permutation is $k$-universal
if it contains every element of $S_k$. Let $\Pi_n$ be uniformly
distributed on $S_n$.

Determine the asymptotic threshold for $k$-universality as
$k\to\infty$. In the standard formulation of Alon's conjecture,
for every fixed $\varepsilon>0$,
\[
 \Pr\!\left(\Pi_{\lceil(1+\varepsilon)k^2/4\rceil}
                 \text{ is }k\text{-universal}\right)\longrightarrow1.
\]
The matching lower-threshold obstruction is the increasing
pattern: below $(1-\varepsilon)k^2/4$, the probability of
universality tends to zero, by longest-increasing-subsequence
asymptotics. The unresolved direction requires all $k!$ patterns
simultaneously, not merely any one prescribed pattern.

Occurrences need not occupy consecutive positions, and their
actual values need only agree in relative order. The parameter
$k$ grows, rather than remaining fixed while $n$ tends to infinity.
The collected Cooper entry writes $k^2/2$; the primary modern
statement by He and Kwan specifies $k^2/4$. This definition records
that explicit constant correction while retaining the source's
question about the universality threshold.
\subsection{Short English statement}
Determine the threshold at which a uniformly random permutation contains every pattern of a growing size k; Alon’s standard conjectured scale is k²/4.
\subsection{Sources}
\begin{itemize}
\item[S1] ULAM.AI and the UnsolvedMath Contributors, supplied problems.json, record 3100083 (AMR-030-0083), AMR Open Problem Lists. Catalogue identity and source transcription; CC BY 4.0.
\url{https://huggingface.co/datasets/ulamai/UnsolvedMath}.
\item[S2] Cooper - Combinatorial Problems I Like (2020 snapshot), Probability, source-order bullet 83. Original source indexed by AMR-030-0083.
\url{https://people.math.sc.edu/cooper/combprob.html}.
\item[S3] He and Kwan, Universality of random permutations, statement of Alon’s conjecture.
\url{https://arxiv.org/abs/1911.12878}.
\end{itemize}
\end{document}
